\documentclass[10pt]{article}
\usepackage[T1]{fontenc}
\usepackage[utf8]{inputenc}
\usepackage[french]{babel}
\usepackage{amsmath,amssymb,mathrsfs,mathpazo}
\usepackage{geometry}
\geometry{a4paper,margin=2cm}
\usepackage{graphicx,array,tabularx,tabularray,colortbl,multirow,multicol,booktabs,enumitem,pifont,ulem,fancybox,diagbox,cancel}
\usepackage{pstricks,pstricks-add,pst-plot,pst-tree,pst-node,pst-text,pst-func,pst-math,pst-bspline,pst-3dplot,pst-eucl,pst-solides3d}
\let\pscalculate\undefined
\usepackage{pst-bezier}
\usepackage{tikz,pgfplots,tkz-tab}
\usetikzlibrary{arrows,arrows.meta,patterns,calc,positioning,trees,decorations.markings,shapes.geometric,shapes.misc,shapes.arrows,intersections,decorations.pathmorphing,decorations.pathreplacing,backgrounds,shadows,babel}
\pgfplotsset{compat=1.18}
\usepgfplotslibrary{fillbetween}
\usepackage[np]{numprint}
\usepackage{listings,makecell,pas-tableur,twoopt}
\let\barre\undefined
\usepackage[pstricks]{tablvar}
\definecolor{darkgreen}{RGB}{0,100,0}
\providecommand{\red}{\color{red}}
\providecommand{\blue}{\color{blue}}
\providecommand{\green}{\color{green}}
\providecommand{\black}{\color{black}}
\providecommand{\white}{\color{white}}
\providecommand{\gray}{\color{gray}}
\providecommand{\cyan}{\color{cyan}}
\providecommand{\magenta}{\color{magenta}}
\providecommand{\yellow}{\color{yellow}}
\providecommand{\violet}{\color{violet}}
\providecommand{\orange}{\color{orange}}
\providecommand{\euro}{\texteuro}
\providecommand{\EUR}{\texteuro}
\providecommand{\ieme}{\textsuperscript{e}}
\providecommand{\degres}{\ensuremath{^\circ}}
\providecommand{\degree}{\ensuremath{^\circ}}
\providecommand{\vv}[1]{\overrightarrow{#1}}
\providecommand{\Coord}[1]{#1}
\providecommand{\decosix}{}
\providecommand{\starredbullet}{\ensuremath{\bullet}}
\pagestyle{empty}
\setlength{\parindent}{0pt}
\usepackage{adjustbox}\DeclareUnicodeCharacter{2236}{:}\DeclareUnicodeCharacter{2190}{\ensuremath{\leftarrow}}
\providecommand{\R}{\mathbb{R}}
\providecommand{\N}{\mathbb{N}}
\providecommand{\Z}{\mathbb{Z}}
\providecommand{\Q}{\mathbb{Q}}
\providecommand{\D}{\mathbb{D}}
\providecommand{\C}{\mathbb{C}}
\providecommand{\e}{{\mathrm{e}}}
\providecommand{\vect}[1]{\overrightarrow{#1}}
\providecommand{\vectt}[1]{\overrightarrow{\mathrm{#1}}}
\providecommand{\vvt}[1]{\overrightarrow{\mathrm{#1}}}
\providecommand{\np}[2][]{{#2~#1}}
\providecommand{\vv}[1]{{\overrightarrow{#1}}}
\renewcommand{\d}{{\mathrm{d}}}
\providecommand{\degres}{{^\circ}}
\providecommand{\degre}{{^\circ}}
\providecommand{\blue}{\color{blue}}
\providecommand{\red}{\color{red}}
\providecommand{\enskip}{\quad}
\providecommand{\llbracket}{\lbrack\!\lbrack}
\providecommand{\rrbracket}{\rbrack\!\rbrack}
\providecommand{\euro}{{\text{€}}}
\providecommand{\no}{{\text{n°}}}
\providecommand{\interval}[2]{{[#1;#2]}}
\providecommand{\intervFO}[2]{{[#1;#2[}}
\providecommand{\dotfill}{\ldots}

\providecommand{\ds}{\displaystyle}
\providecommand{\textding}[1]{#1}
\providecommand{\Oij}{\(\mathrm{O};\vec{\imath},\vec{\jmath}\)}
\providecommand{\Oijk}{\(\mathrm{O};\vec{\imath},\vec{\jmath},\vec{k}\)}
\providecommand{\Ouv}{\(\mathrm{O};\vec{u},\vec{v}\)}
\providecommand{\intervFF}[2]{\ensuremath{[#1;#2]}}
\renewcommand{\euro}{\ifmmode\text{\texteuro}\else\texteuro\fi}
\DeclareUnicodeCharacter{00B0}{\ifmmode^{\circ}\else\textdegree\fi}
\DeclareUnicodeCharacter{20AC}{\ifmmode\text{\texteuro}\else\texteuro\fi}
% Énoncés : Éducation nationale. Transcription et corrigés : APMEP, Denis Vergès et contributeurs. Découpage et classement : CoopMaths, Sébastien Lozano, Éric Elter et contributeurs. Adaptation de la présentation : Maths Explorer.
% CC BY-SA : https://coopmaths.fr/www/about/
% https://coopmaths.fr/alea/static/bac/2023/tex/bac_2023_03_sujet2_reunion_2.tex
\begin{document}
\noindent\textit{APMEP / CoopMaths. CC BY-SA. Présentation adaptée par Maths Explorer.}\par\medskip
\medskip

On considère la suite $\left(u_n\right)$ définie par $u_0 = 8$ et, pour tout entier naturel $n$,
\[u_{n +1} = \dfrac{6u_n + 2}{u_n +5}.\]

\smallskip

\begin{enumerate}
\item Calculer $u_1$.
\item Soit $f$ la fonction définie sur l'intervalle $[0~;~+\infty[$ par :
\[f(x) = \dfrac{6x + 2 }{x +5}.\]

Ainsi, pour tout entier naturel $n$, on a : $u_{n+1} = f\left(u_n\right)$.
	\begin{enumerate}
		\item Démontrer que la fonction $f$ est strictement croissante sur l'intervalle $[0~;~+\infty[$.
		
En déduire que pour tout réel $x > 2$, on a $f(x) > 2$.
		\item Démontrer par récurrence que, pour tout entier naturel n, on a $u_n > 2$.
	\end{enumerate}
\item On admet que, pour tout entier naturel $n$, on a :
\[u_{n+1} - u_n = \dfrac{\left(2 - u_n\right)\left(u_n + 1\right)}{u_n +5}.\]

	\begin{enumerate}
		\item Démontrer que la suite $\left(u_n\right)$ est décroissante.
		\item En déduire que la suite $\left(u_n\right)$ est convergente.
	\end{enumerate}
\item On définit la suite $\left(v_n\right)$ pour tout entier naturel par:
\[v_n = \dfrac{u_n - 2}{u_n + 1}.\]

	\begin{enumerate}
		\item Calculer $v_0$.
		\item Démontrer que $\left(v_n\right)$ est une suite géométrique de raison $\dfrac47$.
		\item Déterminer, en justifiant, la limite de $\left(v_n\right)$.
		
En déduire la limite de $\left(u_n\right)$.
	\end{enumerate}
\end{enumerate}
\begin{minipage}{0.53\linewidth}
\begin{enumerate}[resume]
		\item  On considère la fonction
Python seuil ci-contre, où A est un nombre réel strictement plus grand que 2.

Donner, sans justification, la valeur renvoyée par la commande seuil (2.001) puis interpréter cette valeur dans le contexte de l'exercice.
\end{enumerate}
\end{minipage}\hfill
\begin{minipage}{0.43\linewidth}
\begin{center}
\begin{tabular}{|l l|}\hline
\textbf{def}&seuil (A) :\\
&n = 0\\
&u = 8\\
&\textbf{while} u $>$ A :\\
&\qquad u = (6*u + 2)/(u + 5)\\
&\qquad n = n + 1\\
&\textbf{return} n\\ \hline
\end{tabular}
\end{center}
\end{minipage}

\bigskip

\end{document}
